% ============================================================================
% WORKBOOK — Section 8.2: Making Inferences from Random Samples
% CCSS 7.SP.A.2
% Practice-focused reinforcement for the study guide topic
% ============================================================================
\section{Making Inferences from Random Samples}
\topicTitle{Making Inferences from Random Samples}

\begin{quickReview}{Making Inferences from Random Samples}
An \textbf{inference} is a conclusion about a \emph{population} based on data from a \emph{sample}.

\begin{itemize}[leftmargin=*]
    \item Use the sample proportion or sample mean to \textbf{estimate} the population value.
    \item \textbf{Multiple samples} of the same size may give slightly different results — this is normal variation.
    \item Larger random samples give \textbf{more reliable} estimates.
    \item To estimate a population total: $\text{sample proportion} \times \text{population size}$.
\end{itemize}

\medskip
\textbf{Example:} In a random sample of $40$ students, $10$ prefer art class. Sample proportion $= \frac{10}{40} = 25\%$. If the school has $600$ students, estimate: $0.25 \times 600 = 150$ students prefer art.
\end{quickReview}

% ── Warm-Up ──────────────────────────────────────────────────────────────────
\begin{practiceBox}[funGreen]{\faSun~Warm-Up}

\practiceHeader[funGreen]{Find the Sample Proportion}

\begin{multicols}{2}
\prob $12$ out of $40$ prefer pizza. \answerBlank[3cm]
\answerExplain{$30\%$}{$\frac{12}{40} = 0.30 = 30\%$. The sample proportion is the number with the trait divided by the sample size.}

\prob $9$ out of $30$ ride a bike to school. \answerBlank[3cm]
\answerExplain{$30\%$}{$\frac{9}{30} = 0.30 = 30\%$.}

\prob $20$ out of $80$ chose blue as favorite color. \answerBlank[3cm]
\answerExplain{$25\%$}{$\frac{20}{80} = 0.25 = 25\%$.}

\prob $15$ out of $50$ have a pet dog. \answerBlank[3cm]
\answerExplain{$30\%$}{$\frac{15}{50} = 0.30 = 30\%$.}

\prob $7$ out of $25$ scored above $90$. \answerBlank[3cm]
\answerExplain{$28\%$}{$\frac{7}{25} = 0.28 = 28\%$.}

\prob $18$ out of $60$ prefer reading. \answerBlank[3cm]
\answerExplain{$30\%$}{$\frac{18}{60} = 0.30 = 30\%$.}
\end{multicols}
\end{practiceBox}

% ── Practice A: Estimating Population Totals ─────────────────────────────────
\begin{practiceBox}{Estimating Population Totals}

\practiceHeader{Use sample data to estimate population totals.}

\prob In a random sample of $50$ students, $20$ prefer soccer. The school has $400$ students. Estimate how many prefer soccer. \answerBlank[3cm]
\answerExplain{$160$ students}{Sample proportion: $\frac{20}{50} = 0.40$. Estimate: $0.40 \times 400 = 160$.}

\prob A random sample of $60$ shoppers shows $15$ buy organic produce. The store has $900$ regular shoppers. Estimate the number who buy organic. \answerBlank[3cm]
\answerExplain{$225$ shoppers}{Sample proportion: $\frac{15}{60} = 0.25$. Estimate: $0.25 \times 900 = 225$.}

\prob A survey of $40$ randomly chosen employees found that $28$ use public transit. The company has $500$ employees. Estimate the total number who use public transit. \answerBlank[3cm]
\answerExplain{$350$ employees}{Sample proportion: $\frac{28}{40} = 0.70$. Estimate: $0.70 \times 500 = 350$.}

\prob Out of $25$ randomly tested batteries, $2$ were defective. A shipment has $5{,}000$ batteries. Estimate the number of defective batteries. \answerBlank[3cm]
\answerExplain{$400$ batteries}{Sample proportion: $\frac{2}{25} = 0.08$. Estimate: $0.08 \times 5{,}000 = 400$.}

\prob A random sample of $80$ trees in a forest shows $12$ are oak trees. The forest has $2{,}000$ trees. Estimate the number of oak trees. \answerBlank[3cm]
\answerExplain{$300$ oak trees}{Sample proportion: $\frac{12}{80} = 0.15$. Estimate: $0.15 \times 2{,}000 = 300$.}

\end{practiceBox}

% ── Practice B: Comparing Multiple Samples ───────────────────────────────────
\begin{practiceBox}[funTeal]{Comparing Multiple Samples}

\practiceHeader[funTeal]{Three random samples of $30$ students each were taken from the same school. Use the data to answer questions.}

\begin{center}
\begin{tabular}{lcc}
    \textbf{Sample} & \textbf{Students who prefer math} & \textbf{Sample proportion} \\
    \hline
    Sample 1 & $8$  & \rule{2cm}{0.4pt} \\
    Sample 2 & $10$ & \rule{2cm}{0.4pt} \\
    Sample 3 & $7$  & \rule{2cm}{0.4pt} \\
\end{tabular}
\end{center}

\prob Find the sample proportion for each sample, rounded to the nearest percent.
\answerExplain{Sample 1: $27\%$; Sample 2: $33\%$; Sample 3: $23\%$}{Sample 1: $\frac{8}{30} \approx 0.267 \approx 27\%$. Sample 2: $\frac{10}{30} \approx 0.333 \approx 33\%$. Sample 3: $\frac{7}{30} \approx 0.233 \approx 23\%$.}

\prob Find the average of the three sample proportions. \answerBlank[3cm]
\answerExplain{$\approx 28\%$}{$\frac{27 + 33 + 23}{3} \approx 27.7\% \approx 28\%$. Averaging multiple samples gives a more reliable estimate.}

\prob If the school has $450$ students, use the average proportion to estimate how many prefer math. \answerBlank[3cm]
\answerExplain{$\approx 125$ students}{$0.28 \times 450 \approx 126$. Rounding, about $125$ students prefer math.}

\prob Why do the three samples give different proportions?
\answerExplain{Random variation}{Each sample contains different students, so the proportions vary. This is normal and expected with random sampling.}

\end{practiceBox}

% ── Practice C: True or False ────────────────────────────────────────────────
\begin{practiceBox}[funPurple]{True or False?}

\trueOrFalse{A sample mean is always exactly equal to the population mean.}
\answerExplain{False}{A sample mean is an estimate of the population mean. Due to random variation, it is rarely exactly equal but should be close if the sample is random and large enough.}

\trueOrFalse{Taking multiple random samples and averaging the results gives a better estimate than using a single sample.}
\answerExplain{True}{Averaging multiple samples reduces the effect of random variation, giving a more reliable estimate of the population value.}

\trueOrFalse{If a sample of $100$ people shows $40\%$ like chocolate ice cream, then exactly $40\%$ of the population likes chocolate ice cream.}
\answerExplain{False}{The $40\%$ is an estimate based on the sample. The actual population proportion is likely close to $40\%$ but not necessarily exactly $40\%$.}

\trueOrFalse{A random sample of $500$ people is generally more reliable than a random sample of $50$ people.}
\answerExplain{True}{Larger samples reduce variability and produce estimates closer to the true population value.}

\end{practiceBox}

% ── Word Problems ────────────────────────────────────────────────────────────
\begin{practiceBox}[funBlue]{\faGlobe~Word Problems}

\wordProblem{A bakery randomly samples $30$ cookies from a batch of $600$ and finds $3$ are undercooked. Estimate the total number of undercooked cookies in the batch.}{cookies}
\answerExplain{$60$ cookies}{Sample proportion: $\frac{3}{30} = 0.10 = 10\%$. Estimate: $0.10 \times 600 = 60$ undercooked cookies.}

\wordProblem{A random survey of $75$ townspeople finds that $45$ support building a new library. The town has $6{,}000$ residents. About how many support the library?}{residents}
\answerExplain{$3{,}600$ residents}{Sample proportion: $\frac{45}{75} = 0.60 = 60\%$. Estimate: $0.60 \times 6{,}000 = 3{,}600$.}

\wordProblem{Two random samples of fish in a lake are collected. Sample 1: $8$ out of $40$ are trout. Sample 2: $12$ out of $40$ are trout. The lake has an estimated $2{,}000$ fish. Use the average proportion to estimate the number of trout.}{trout}
\answerExplain{$500$ trout}{Sample 1 proportion: $\frac{8}{40} = 0.20$. Sample 2: $\frac{12}{40} = 0.30$. Average: $\frac{0.20 + 0.30}{2} = 0.25$. Estimate: $0.25 \times 2{,}000 = 500$ trout.}

\end{practiceBox}

% ── Challenge ────────────────────────────────────────────────────────────────
\begin{challengeBox}
\prob A school of $1{,}200$ students collects three random samples of $50$ students each. The number who eat breakfast every day is $30$, $28$, and $32$. Use all three samples to estimate the total number of students who eat breakfast daily. Show your work. \answerBlank[3cm]
\answerExplain{$720$ students}{Average number per sample: $\frac{30+28+32}{3} = 30$. Average proportion: $\frac{30}{50} = 0.60$. Estimate: $0.60 \times 1{,}200 = 720$ students.}

\prob Sarah surveys $40$ random students and finds $25\%$ want a longer recess. Jake surveys $100$ random students and finds $30\%$ want a longer recess. Whose estimate is likely more reliable, and why?
\answerExplain{Jake's estimate}{Jake's sample is larger ($100$ vs. $40$), which reduces the effect of random variation. Larger samples tend to produce estimates closer to the true population value.}
\end{challengeBox}

\encouragement{Great job making predictions from data! You are thinking like a real statistician!}
